% ===================================================================
%  LES SIGNETS DU PDF — la table des matieres logicielle.
%  Ecrits avec la primitive \pdfoutline de pdfTeX, et NON avec hyperref :
%  hyperref redefinit \parindent, \baselineskip et la mise en boite des
%  flottants, et ce volume a sa pagination gelee au dixieme de
%  millimetre. La primitive, elle, n'ecrit que dans le catalogue du PDF
%  et ne touche a rien de la composition. Verifie : 236 pages avant,
%  236 apres, aucune ligne deplacee.
%  « goto page N » vise la page par son rang, sans ancre a poser dans le
%  texte — donc sans le moindre objet ajoute a la page.
%  « count -n » ouvre une entree fermee portant n enfants.
%  Le texte des signets est reduit a l'ASCII : le catalogue PDF
%  n'accepte que PDFDocEncoding ou l'UTF-16, et les titres du volume
%  n'ont d'accents qu'aux chevrons, remplaces par des guillemets droits.
% ===================================================================
\pdfoutline goto page 1 {/Fit} count 0 {Kovrilo}
\pdfoutline goto page 5 {/Fit} count 0 {Titulo}
\pdfoutline goto page 6 {/Fit} count 0 {Konstato}
\pdfoutline goto page 7 {/Fit} count 0 {Averto}
\pdfoutline goto page 11 {/Fit} count -23 {Unesma parto - MORFOLOGIO E SINTAXO}
  \pdfoutline goto page 13 {/Fit} count 0 {ALFABETO}
  \pdfoutline goto page 13 {/Fit} count 0 {PRONUNCO DIL VOKALI}
  \pdfoutline goto page 14 {/Fit} count 0 {PRONUNCO DIL KONSONANTI E DIGRAMI}
  \pdfoutline goto page 17 {/Fit} count 0 {ACENTO TONIKA}
  \pdfoutline goto page 20 {/Fit} count 0 {ARTIKLO}
  \pdfoutline goto page 23 {/Fit} count 0 {SUBSTANTIVO}
  \pdfoutline goto page 26 {/Fit} count 0 {PROPRA NOMI}
  \pdfoutline goto page 29 {/Fit} count 0 {ADJEKTIVO QUALIFIKANTA}
  \pdfoutline goto page 33 {/Fit} count 0 {GRADI KOMPARALA}
  \pdfoutline goto page 34 {/Fit} count 0 {PERSONAL PRONOMI}
  \pdfoutline goto page 36 {/Fit} count 0 {POSEDAL ADJEKTIVI E PRONOMI}
  \pdfoutline goto page 38 {/Fit} count 0 {DEMONSTRATIV ADJEKTIVI-PRONOMI}
  \pdfoutline goto page 40 {/Fit} count 0 {RELATIVA E QUESTIONALA ADJEKTIVI-PRONOMI}
  \pdfoutline goto page 41 {/Fit} count 0 {PRONOMO " LO "}
  \pdfoutline goto page 42 {/Fit} count 0 {ADJEKTIVI-PRONOMI NEDEFINITA}
  \pdfoutline goto page 48 {/Fit} count 0 {VERBO}
  \pdfoutline goto page 58 {/Fit} count 0 {ADVERBI}
  \pdfoutline goto page 69 {/Fit} count 0 {PREPOZICIONI}
  \pdfoutline goto page 90 {/Fit} count 0 {KONJUNCIONI}
  \pdfoutline goto page 95 {/Fit} count 0 {INTERJECIONI}
  \pdfoutline goto page 96 {/Fit} count 0 {NOMBRI}
  \pdfoutline goto page 102 {/Fit} count 0 {SINTAXO}
  \pdfoutline goto page 108 {/Fit} count 0 {TEMPI E MODI}
\pdfoutline goto page 117 {/Fit} count -24 {Duesma parto - VORTIFADO}
  \pdfoutline goto page 119 {/Fit} count 0 {ELEMENTI DI VORTO}
  \pdfoutline goto page 119 {/Fit} count 0 {PROCEDI DI VORTIFADO}
  \pdfoutline goto page 120 {/Fit} count 0 {REGULI DI DERIVADO *}
  \pdfoutline goto page 123 {/Fit} count 0 {DERIVADO PER AFIXI}
  \pdfoutline goto page 124 {/Fit} count 0 {AFIXI}
  \pdfoutline goto page 124 {/Fit} count 0 {PREFIXI}
  \pdfoutline goto page 132 {/Fit} count 0 {PREFIXI TEKNIKALA}
  \pdfoutline goto page 133 {/Fit} count 0 {REMARKO PRI LA AFIXI}
  \pdfoutline goto page 134 {/Fit} count 0 {PREPOZICIONI PREFIXA}
  \pdfoutline goto page 134 {/Fit} count 0 {SUFIXI}
  \pdfoutline goto page 170 {/Fit} count 0 {KOMPOZADO}
  \pdfoutline goto page 171 {/Fit} count 0 {REGULO DI ANALIZO O DESKOMPOZO}
  \pdfoutline goto page 174 {/Fit} count 0 {KOMPOZADO PER PREPOZICIONI}
  \pdfoutline goto page 175 {/Fit} count 0 {LA KOMPOZAJI E LA SUFIXI}
  \pdfoutline goto page 181 {/Fit} count 0 {L'ACENTIZO EN IDO}
  \pdfoutline goto page 185 {/Fit} count 0 {LA PLURALO EN IDO}
  \pdfoutline goto page 190 {/Fit} count 0 {GENRO E MASKULISMO}
  \pdfoutline goto page 197 {/Fit} count 0 {SUBSTANTIVIGO DIL ADJEKTIVO}
  \pdfoutline goto page 207 {/Fit} count 0 {CA, TA e QUA}
  \pdfoutline goto page 208 {/Fit} count 0 {KONJUGO-SISTEMO DI IDO}
  \pdfoutline goto page 217 {/Fit} count 0 {VORTORDINO}
  \pdfoutline goto page 219 {/Fit} count 0 {PUNTIZADO}
  \pdfoutline goto page 224 {/Fit} count 0 {NOMI. ADRESI}
  \pdfoutline goto page 225 {/Fit} count 0 {FORMULI DI POLITESO EN LETRI}
\pdfoutline goto page 226 {/Fit} count 0 {TABELO DI LA MATERIO}
\pdfoutline goto page 234 {/Fit} count 0 {APENDICI}
\pdfoutline goto page 234 {/Fit} count 0 {MONDO-BIBLIOTEKO}
\pdfoutline goto page 234 {/Fit} count 0 {RADIO-LEXIKI}
